\chapter{Sobolev Spaces and Weak Derivatives}\label{sec:sobolev}
%% ========================================================================

% Section 2: Sobolev Spaces and Weak Derivatives
% Translation of:
%   DeGiorgi/SobolevSpace.lean
%   DeGiorgi/SobolevSpace/WeakDerivatives.lean
%   DeGiorgi/SobolevSpace/Witnesses.lean
%   DeGiorgi/SobolevSpace/Approximation.lean
%   DeGiorgi/SobolevSpace/PositivePartPrelude.lean
%   DeGiorgi/WholeSpaceSobolev.lean

Throughout this section we fix a positive integer $d$ and write $E = \R^d$ (in
the Lean formalization, $E$ is the Euclidean space of dimension $d$, with the
canonical orthonormal basis $(e_i)_{i=1}^d$). All open sets $\Omega \subseteq E$
are equipped with Lebesgue measure. For a smooth scalar test function $\varphi
\in C_c^\infty(E)$, $i \in \{1,\ldots,d\}$ and $x \in E$ we abbreviate the
$i$-th partial derivative by $\partial_i \varphi(x) = (D\varphi(x))(e_i)$.

\section{Weak partial derivatives and weak gradients}

\begin{definition}
\label{def:HasWeakPartialDeriv}
\lean{DeGiorgi.HasWeakPartialDeriv}
\leanok
Let $\Omega \subseteq E$ be a set, $i \in \{1,\ldots,d\}$, and $f, g : E \to \R$.
We say that $g$ is the \emph{weak $i$-th partial derivative of $f$ on $\Omega$},
written $\partial_i^w f = g$ on $\Omega$, if for every test function
$\varphi \in C^\infty(E,\R)$ with compact support and $\supp \varphi \subseteq
\Omega$ we have
\begin{equation*}
  \int_\Omega f(x) \, \partial_i \varphi(x) \, dx
    = - \int_\Omega g(x) \, \varphi(x) \, dx.
\end{equation*}
\end{definition}

\begin{definition}
\label{def:HasWeakGrad}
\lean{DeGiorgi.HasWeakGrad, DeGiorgi.HasWeakDiv}
\leanok
\uses{def:HasWeakPartialDeriv}
A vector field $G : E \to E$ is a \emph{weak gradient} of $f$ on $\Omega$,
written $\nabla^w f = G$, if each component $G_i$ is the weak $i$-th partial
derivative of $f$ on $\Omega$.

A scalar function $g : E \to \R$ is a \emph{weak divergence} of a vector field
$F : E \to E$ on $\Omega$ if for every $\varphi \in C_c^\infty$ with $\supp
\varphi \subseteq \Omega$
\begin{equation*}
  \int_\Omega \sum_{i=1}^d F_i(x) \, \partial_i \varphi(x) \, dx
    = - \int_\Omega g(x)\, \varphi(x)\, dx.
\end{equation*}
\end{definition}

\begin{lemma}[Uniqueness]
\label{lem:weak_partial_unique}
\lean{DeGiorgi.HasWeakPartialDeriv.ae_eq}
\leanok
\uses{def:HasWeakPartialDeriv}
Let $\Omega \subseteq E$ be open, $i \in \{1,\ldots,d\}$ and $f, g_1, g_2 :
E \to \R$ with $g_1$ and $g_2$ locally integrable on $\Omega$. If both $g_1$
and $g_2$ are weak $i$-th partial derivatives of $f$ on $\Omega$, then
$g_1 = g_2$ almost everywhere on $\Omega$.
\end{lemma}

\begin{proof}
\leanok

It suffices to show that $g_1 - g_2 = 0$ a.e.\ on $\Omega$. The function $g_1 -
g_2$ is locally integrable on $\Omega$, and for any $\varphi \in C_c^\infty(E)$
with $\supp \varphi \subseteq \Omega$ we have, subtracting the two test
identities,
\begin{equation*}
  \int_\Omega g_1(x)\,\varphi(x)\,dx = \int_\Omega g_2(x)\,\varphi(x)\,dx,
\end{equation*}
so $\int_\Omega (g_1 - g_2)(x)\,\varphi(x)\,dx = 0$. Since $\varphi$ vanishes
outside its support which lies in $\Omega$, this is the same as integrating over
the whole space $E$. By the standard fact that a locally integrable function
whose integral against every smooth compactly supported test function vanishes
is zero almost everywhere (\leanname{IsOpen.ae\_eq\_zero\_of\_integral\_contDiff\_smul\_eq\_zero}),
we conclude $g_1 = g_2$ a.e.\ on $\Omega$.
\end{proof}

\begin{lemma}[Classical $\Rightarrow$ weak derivative]
\label{lem:weak_of_contDiff}
\lean{DeGiorgi.HasWeakPartialDeriv.of_contDiff}
\leanok
\uses{def:HasWeakPartialDeriv}
Let $\Omega \subseteq E$ be open and $f \in C^1(E,\R)$. Then $\partial_i f$ is
the weak $i$-th partial derivative of $f$ on $\Omega$.
\end{lemma}

\begin{proof}
\leanok

Take any $\varphi \in C_c^\infty(E)$ with $\supp\varphi \subseteq \Omega$. Both
$f \cdot \partial_i \varphi$ and $\partial_i f \cdot \varphi$ vanish outside
$\supp\varphi \subseteq \Omega$, so the corresponding integrals on $\Omega$ and
on $E$ coincide. The classical integration-by-parts formula
(\leanname{integral\_mul\_fderiv\_eq\_neg\_fderiv\_mul\_of\_integrable}) for
differentiable functions with compactly supported product gives
\begin{equation*}
  \int_E f(x)\,\partial_i \varphi(x)\,dx
    = - \int_E (\partial_i f)(x)\,\varphi(x)\,dx,
\end{equation*}
which is the weak-derivative identity on $\Omega$.
\end{proof}

\begin{lemma}[Product rule against a smooth factor]
\label{lem:weak_product_rule}
\lean{DeGiorgi.HasWeakPartialDeriv.mul_smooth}
\leanok
\uses{def:HasWeakPartialDeriv}
Let $\Omega \subseteq E$ be open, $i \in \{1,\ldots,d\}$, and let $f, g, \eta :
E \to \R$ with $\eta \in C^\infty$, $f$ and $g$ locally integrable on $\Omega$,
and $g$ a weak $i$-th partial derivative of $f$ on $\Omega$. Then
\begin{equation*}
  x \mapsto \eta(x)\,g(x) + (\partial_i \eta(x))\,f(x)
\end{equation*}
is the weak $i$-th partial derivative of $\eta f$ on $\Omega$.
\end{lemma}

\begin{proof}
\leanok

Fix a test function $\varphi \in C_c^\infty(E)$ with $\supp\varphi \subseteq
\Omega$. Then $\eta \varphi$ is also smooth, compactly supported, with support
contained in $\Omega$. Applying the weak-derivative identity for $f$ to the test
function $\eta \varphi$ and using the smooth product rule
$\partial_i(\eta\varphi) = \eta\,\partial_i\varphi + \varphi\,\partial_i\eta$
gives
\begin{equation*}
  \int_\Omega f(x)\bigl(\eta(x)\,\partial_i\varphi(x) + \varphi(x)\,\partial_i\eta(x)\bigr)\,dx
    = -\int_\Omega g(x)\,\eta(x)\,\varphi(x)\,dx.
\end{equation*}
The local integrability of $f$ and $g$, together with continuity of $\eta$,
$\partial_i\eta$, $\varphi$, $\partial_i\varphi$ and compactness of supports,
ensures that all integrals are absolutely convergent and may be split. Splitting
the left-hand side and rearranging yields
\begin{align*}
  \int_\Omega \eta(x)\,f(x)\,\partial_i\varphi(x)\,dx
  &= -\int_\Omega g(x)\,\eta(x)\,\varphi(x)\,dx
     - \int_\Omega f(x)\,\varphi(x)\,\partial_i\eta(x)\,dx \\
  &= -\int_\Omega \bigl(\eta(x)g(x) + (\partial_i\eta(x))f(x)\bigr)\,\varphi(x)\,dx,
\end{align*}
which is the desired weak-derivative identity for $\eta f$.
\end{proof}

\begin{lemma}[Restriction to a sub-open set]
\label{lem:weak_restrict}
\lean{DeGiorgi.HasWeakPartialDeriv.restrict}
\leanok
\uses{def:HasWeakPartialDeriv}
Let $\Omega' \subseteq \Omega \subseteq E$ with $\Omega'$ open. If $g$ is the
weak $i$-th partial derivative of $f$ on $\Omega$, then it is also the weak
$i$-th partial derivative of $f$ on $\Omega'$.
\end{lemma}

\begin{proof}
\leanok

Let $\varphi \in C_c^\infty(E)$ with $\supp\varphi \subseteq \Omega'$. Then a
fortiori $\supp\varphi \subseteq \Omega$, so applying the hypothesis on
$\Omega$ gives the integral identity over $\Omega$. Outside $\Omega'$ the test
function $\varphi$ and the field $\partial_i\varphi$ both vanish (the latter
because $\varphi$ is locally constant equal to zero), so the integrals over
$\Omega$ and $\Omega'$ of $f\,\partial_i\varphi$ and $g\,\varphi$ agree. The
result follows.
\end{proof}

\begin{lemma}[$L^p$-closure of weak partial derivatives]
\label{lem:weak_lp_closure}
\lean{DeGiorgi.HasWeakPartialDeriv.of_eLpNormApprox_p}
\leanok
\uses{def:HasWeakPartialDeriv}
Let $\Omega \subseteq E$ be open and $1 < p < \infty$. Suppose $f, g \in
\Lp{p}(\Omega)$ and there are sequences $(\psi_n)$ and $(g_{\psi_n})$ such that:
\begin{itemize}
  \item $g_{\psi_n}$ is the weak $i$-th partial derivative of $\psi_n$ on
        $\Omega$ for every $n$;
  \item $\psi_n - f$ and $g_{\psi_n} - g$ lie in $\Lp{p}(\Omega)$ and tend to
        $0$ in $\Lp{p}(\Omega)$ as $n \to \infty$.
\end{itemize}
Then $g$ is the weak $i$-th partial derivative of $f$ on $\Omega$.
\end{lemma}

\begin{proof}
\leanok

Fix $\varphi \in C_c^\infty(E)$ with $\supp\varphi \subseteq \Omega$. Let $q$
denote the H{\"o}lder conjugate of $p$, so that $1 < q < \infty$ and $1/p + 1/q
= 1$. Both $\varphi$ and $\partial_i\varphi$ are continuous, compactly
supported, and hence belong to $\Lp{q}(\Omega)$.

By H{\"o}lder's inequality on $(\Omega, dx)$,
\begin{align*}
  \left| \int_\Omega \partial_i\varphi(x) (\psi_n(x) - f(x))\,dx \right|
  &\le \|\partial_i\varphi\|_{\Lp{q}(\Omega)}\,\|\psi_n - f\|_{\Lp{p}(\Omega)},\\
  \left| \int_\Omega \varphi(x) (g_{\psi_n}(x) - g(x))\,dx \right|
  &\le \|\varphi\|_{\Lp{q}(\Omega)}\,\|g_{\psi_n} - g\|_{\Lp{p}(\Omega)},
\end{align*}
both of which tend to zero by hypothesis. Hence
\begin{align*}
  \int_\Omega \psi_n(x)\,\partial_i\varphi(x)\,dx
    &\;\to\; \int_\Omega f(x)\,\partial_i\varphi(x)\,dx, \\
  \int_\Omega g_{\psi_n}(x)\,\varphi(x)\,dx
    &\;\to\; \int_\Omega g(x)\,\varphi(x)\,dx.
\end{align*}
By assumption, for every $n$
\begin{equation*}
  \int_\Omega \psi_n(x)\,\partial_i\varphi(x)\,dx
    = -\int_\Omega g_{\psi_n}(x)\,\varphi(x)\,dx,
\end{equation*}
and passing to the limit yields
\begin{equation*}
  \int_\Omega f(x)\,\partial_i\varphi(x)\,dx
    = -\int_\Omega g(x)\,\varphi(x)\,dx,
\end{equation*}
which is the weak-derivative identity for $f$ and $g$.
\end{proof}

\section{Sobolev membership and weak-gradient witnesses}

\begin{definition}
\label{def:MemW1p}
\lean{DeGiorgi.MemW1p}
\leanok
\uses{def:HasWeakPartialDeriv}
Let $p \in [1,\infty]$, $\Omega \subseteq E$ a set, and $f : E \to \R$. We say
that $f \in \Wkp{1}{p}(\Omega)$ if $f \in \Lp{p}(\Omega)$ and for every $i \in
\{1,\ldots,d\}$ there exists $g_i \in \Lp{p}(\Omega)$ such that $g_i$ is the
weak $i$-th partial derivative of $f$ on $\Omega$.
\end{definition}

\begin{definition}
\label{def:MemW1pWitness}
\lean{DeGiorgi.MemW1pWitness}
\leanok
\uses{def:HasWeakGrad}
A \emph{$\Wkp{1}{p}$-witness} for $f$ on $\Omega$ is a tuple consisting of a
proof that $f \in \Lp{p}(\Omega)$, a vector field $G : E \to E$ (the weak
gradient field) such that each component $x \mapsto G_i(x)$ lies in
$\Lp{p}(\Omega)$, together with a proof that $G$ is a weak gradient of $f$ on
$\Omega$ (Definition~\ref{def:HasWeakGrad}).
\end{definition}

The Hilbert-space case $p = 2$ is also denoted $H^1(\Omega) := W^{1,2}(\Omega)$
in the formalization (\leanname{DeGiorgi.MemH1}).

\begin{definition}
\label{def:MemW01p}
\lean{DeGiorgi.MemW01p}
\leanok
\uses{def:MemW1p, def:MemW1pWitness}
We say that $f \in \Wop{p}(\Omega)$ if $f \in \Wkp{1}{p}(\Omega)$ and there
exist a $\Wkp{1}{p}$-witness $(G)$ for $f$ on $\Omega$ and a sequence
$(\varphi_n)_{n\in\N}$ of smooth compactly supported functions with
$\supp\varphi_n \subseteq \Omega$ such that
\begin{equation*}
  \|\varphi_n - f\|_{\Lp{p}(\Omega)} \to 0
  \quad\text{and}\quad
  \|\partial_i \varphi_n - G_i\|_{\Lp{p}(\Omega)} \to 0
  \;\text{ for every } i.
\end{equation*}
We write $H_0^1(\Omega) := W_0^{1,2}(\Omega)$ (\leanname{DeGiorgi.MemH01}).
\end{definition}

\begin{remark}[Existence of an explicit witness; \leanname{DeGiorgi.MemW1p.someWitness}]
Given $f \in W^{1,p}(\Omega)$, axiom of choice produces a witness as in
Definition~\ref{def:MemW1pWitness}: choose for each coordinate $i$ a function
$g_i$ as in Definition~\ref{def:MemW1p}, and assemble them into the vector
field $G(x) := (g_1(x),\ldots,g_d(x))$.
\end{remark}

\begin{lemma}[Forgetting the witness]
\label{lem:forgetting-witness}
\lean{DeGiorgi.MemW1pWitness.memW1p, DeGiorgi.MemW01p.memW1p}
\leanok
\uses{def:MemW1p, def:MemW1pWitness, def:MemW01p}
(\leanname{DeGiorgi.MemW1pWitness.memW1p},
\leanname{DeGiorgi.MemW01p.memW1p})\\
Every witness gives rise to membership: a $\Wkp{1}{p}$-witness yields
$f \in W^{1,p}(\Omega)$, and an element of $\Wop{p}(\Omega)$ is in particular
in $\Wkp{1}{p}(\Omega)$.
\end{lemma}

\begin{proof}
\leanok

Direct from the definitions: take $g_i := G_i$.
\end{proof}

\begin{lemma}[Sum of witnesses]
\label{lem:witness_add}
\lean{DeGiorgi.MemW1pWitness.add}
\leanok
\uses{def:MemW1pWitness}
For $u, v \in H^1(\Omega)$ with witnesses $G_u$ and $G_v$, the function $u + v$
admits the witness $G_u + G_v$.
\end{lemma}

\begin{proof}
\leanok

Membership of $u+v$ in $\Lp{2}(\Omega)$ and of $G_u + G_v$ in $\Lp{2}(\Omega)$
in each coordinate follows from $\Lp{2}$ being a vector space. To verify the
weak-gradient identity, take $\varphi \in C_c^\infty(E)$ with $\supp\varphi
\subseteq \Omega$. The witnesses $G_u$ and $G_v$ give
\begin{align*}
  \int_\Omega u(x)\,\partial_i\varphi(x)\,dx
    &= -\int_\Omega (G_u)_i(x)\,\varphi(x)\,dx, \\
  \int_\Omega v(x)\,\partial_i\varphi(x)\,dx
    &= -\int_\Omega (G_v)_i(x)\,\varphi(x)\,dx.
\end{align*}
Adding these identities, using local integrability and the linearity of the
integral, yields the weak-gradient identity for $u + v$ with weak gradient
$G_u + G_v$.
\end{proof}

\begin{lemma}[Restriction of witnesses]
\label{lem:MemW1pWitness_restrict}
\lean{DeGiorgi.MemW1pWitness.restrict}
\leanok
\uses{def:MemW1pWitness}
Let $\Omega' \subseteq \Omega$ with $\Omega'$ open. Any
$\Wkp{1}{p}$-witness on $\Omega$ restricts to a $\Wkp{1}{p}$-witness on
$\Omega'$, with the same gradient field.
\end{lemma}

\begin{proof}
\leanok
\uses{lem:weak_restrict}

$\Lp{p}$ membership is monotone in the measure (here the volume measure
restricted to a smaller set), and Lemma~\ref{lem:weak_restrict} gives the weak
gradient identity on $\Omega'$.
\end{proof}

\begin{lemma}[Scalar multiple of a witness]
\label{lem:MemW1pWitness_smul}
\lean{DeGiorgi.MemW1pWitness.smul}
\leanok
\uses{def:MemW1pWitness}
For $u \in H^1(\Omega)$ with witness $G$ and $c \in \R$, the function $cu$
admits the witness $cG$.
\end{lemma}

\begin{proof}
\leanok

The $\Lp{2}$ memberships are immediate. For the test-function identity, fix
$\varphi$ as before; using linearity and the witness for $u$,
\begin{equation*}
  \int_\Omega (cu)(x)\partial_i\varphi(x)\,dx
    = c\int_\Omega u(x)\partial_i\varphi(x)\,dx
    = -c\int_\Omega G_i(x)\varphi(x)\,dx
    = -\int_\Omega (cG)_i(x)\varphi(x)\,dx. \qedhere
\end{equation*}
\end{proof}

\begin{proposition}[Multiplying by a bounded smooth scalar]
\label{prop:witness_mul_smooth_bounded}
\lean{DeGiorgi.MemW1pWitness.mul_smooth_bounded_p}
\leanok
\uses{def:MemW1pWitness}
Let $1 \le p \le \infty$, $\Omega$ open, $u$ a function with $\Wkp{1}{p}$-witness
$G$ on $\Omega$, and $\eta \in C^\infty(E)$ satisfying
$|\eta(x)| \le C_0$ and $\|D\eta(x)\| \le C_1$ for all $x$ and constants
$C_0, C_1 \ge 0$. Then $\eta u$ admits the $\Wkp{1}{p}$-witness whose gradient
field is
\begin{equation*}
  x \mapsto \eta(x)\,G(x) + \bigl(\partial_1\eta(x)\,u(x),\ldots,\partial_d\eta(x)\,u(x)\bigr).
\end{equation*}
\end{proposition}

\begin{proof}
\leanok
\uses{lem:weak_product_rule}

The $\Lp{p}$ bound on $\eta u$ follows from the pointwise inequality
$|\eta(x)u(x)| \le C_0 |u(x)|$ and $u \in \Lp{p}(\Omega)$. For each component
of the displayed gradient field:
\begin{itemize}
  \item $|\eta(x)G_i(x)| \le C_0 |G_i(x)|$, and $G_i \in \Lp{p}(\Omega)$;
  \item $|\partial_i\eta(x)\,u(x)| \le \|D\eta(x)\|\,|u(x)| \le C_1 |u(x)|$,
        and $u \in \Lp{p}(\Omega)$.
\end{itemize}
By Minkowski, the sum lies in $\Lp{p}(\Omega)$. The weak-gradient identity is
the product rule (Lemma~\ref{lem:weak_product_rule}) applied componentwise to
the weak partial derivative $G_i$ of $u$ and the smooth factor $\eta$.
\end{proof}

\begin{lemma}[Smooth compactly supported functions are $\Wkp{1}{p}$]
\label{lem:witness_smooth_cpt}
\lean{DeGiorgi.MemW1pWitness.of_contDiff_hasCompactSupport}
\leanok
\uses{def:MemW1pWitness}
(\leanname{DeGiorgi.MemW1pWitness.of\_contDiff\_hasCompactSupport})\\
Let $f \in C^\infty_c(E,\R)$. Then $f$ admits a canonical $\Wkp{1}{p}$-witness
on the whole space $E$, with weak gradient $\nabla^w f = \nabla f$ given by the
classical gradient field.
\end{lemma}

\begin{proof}
\leanok
\uses{lem:weak_of_contDiff}

Smooth compactly supported functions are integrable to any power, so $f \in
\Lp{p}(E)$ and likewise each classical partial derivative. By
Lemma~\ref{lem:weak_of_contDiff} the classical partial derivatives are weak
partial derivatives.
\end{proof}

\begin{lemma}[Norm of weak gradient lies in $\Lp{p}$]
\label{lem:weak_grad_norm_memLp}
\lean{DeGiorgi.MemW1pWitness.weakGrad_memLp, DeGiorgi.MemW1pWitness.weakGrad_norm_memLp}
\leanok
\uses{def:MemW1pWitness}
Let $G$ be the weak gradient field of a $\Wkp{1}{p}$-witness on $\Omega$. Then
$G \in \Lp{p}(\Omega; E)$, and in particular $\|G(\cdot)\| \in \Lp{p}(\Omega)$.
\end{lemma}

\begin{proof}
\leanok

Each component $G_i$ lies in $\Lp{p}(\Omega)$ by definition of a witness. For
finite-dimensional Euclidean targets the $\Lp{p}$ membership of a vector field
is equivalent to the $\Lp{p}$ membership of each component
(\leanname{MemLp.of\_eval\_piLp}), so $G \in \Lp{p}(\Omega; E)$. Taking norms
preserves $\Lp{p}$ membership.
\end{proof}

\begin{proposition}[$H_0^1$ is a vector space]
\label{prop:H01_vector_space}
\lean{DeGiorgi.MemW01p.add, DeGiorgi.MemW01p.smul, DeGiorgi.MemW01p.sub}
\leanok
\uses{def:MemW01p}
The space $H_0^1(\Omega)$ is closed under addition, scalar multiplication, and
subtraction.
\end{proposition}

\begin{proof}
\leanok
\uses{lem:witness_add, lem:MemW1pWitness_smul}

For addition: Lemma~\ref{lem:witness_add} provides a witness for the sum.
Approximating sequences add componentwise; by the triangle inequality for
$\Lp{2}$ on differences,
\begin{equation*}
  \|\varphi_n^u + \varphi_n^v - (u+v)\|_{\Lp{2}(\Omega)}
  \le \|\varphi_n^u - u\|_{\Lp{2}(\Omega)} + \|\varphi_n^v - v\|_{\Lp{2}(\Omega)}
  \;\to\; 0,
\end{equation*}
and similarly for the gradients (using the linearity of the classical
derivative on smooth functions).

For scalar multiplication by $c$: rescale the approximating sequence by $c$ and
use the homogeneity of the $\Lp{2}$ norm,
$\|c\varphi_n - c u\|_{\Lp{2}} = |c|\,\|\varphi_n - u\|_{\Lp{2}}$.

Subtraction follows by writing $u - v = u + (-1)v$.
\end{proof}

\begin{lemma}[Smooth compactly supported $\Rightarrow$ $\Wop{p}$]
\label{lem:memW01p_of_smooth_cpt}
\lean{DeGiorgi.memW01p_of_contDiff_hasCompactSupport, DeGiorgi.memW01p_of_contDiff_hasCompactSupport_subset}
\leanok
\uses{def:MemW01p}
If $f \in C_c^\infty(E)$ then $f \in \Wop{p}(E)$. More generally, if $\Omega
\subseteq E$ is open and $f \in C_c^\infty(E)$ with $\supp f \subseteq \Omega$,
then $f \in \Wop{p}(\Omega)$.
\end{lemma}

\begin{proof}
\leanok
\uses{lem:weak_of_contDiff, def:MemW1pWitness, lem:witness_smooth_cpt}

Take the constant approximating sequence $\varphi_n := f$. The $\Lp{p}$
differences are identically zero, so trivially converge to $0$, and
Lemma~\ref{lem:witness_smooth_cpt} provides the witness.
\end{proof}

\section{Cutoffs and zero-extension}

\begin{lemma}[Smooth cutoff between a compact set and an open neighborhood]
\label{lem:smooth_cutoff}
\leanok
\noindent\textit{(Formalized as private declaration \leanname{DeGiorgi.exists\_smooth\_cutoff}.)}\par
Let $K \subseteq \Omega \subseteq E$ with $K$ compact and $\Omega$ open. Then
there exists $\eta \in C^\infty(E)$ with compact support such that
$0 \le \eta \le 1$, $\eta = 1$ on $K$, and $\supp\eta \subseteq \Omega$.
\end{lemma}

\begin{proof}
\leanok

Since $K$ is compact and contained in the open set $\Omega$, there exists
$\delta > 0$ such that the closed $\delta$-thickening of $K$ is contained in
$\Omega$. The classical existence of a smooth bump (in Lean,
\leanname{exists\_contMDiff\_support\_eq\_eq\_one\_iff}) yields a smooth $\eta$
with $\eta(x) = 1$ on $K$, support equal to the open $\delta$-thickening of $K$,
and range in $[0,1]$. Its support is then contained in the closed
$\delta$-thickening of $K$ and hence in $\Omega$, and it is automatically
compactly supported.
\end{proof}

\begin{lemma}[Vanishing outside the support]
\label{lem:zero_outside_tsupport}
\lean{DeGiorgi.zero_outside_of_tsupport_subset, DeGiorgi.fderiv_apply_zero_outside_of_tsupport_subset}
\leanok
If $\supp f \subseteq \Omega$ and $x \notin \Omega$, then $f(x) = 0$. If
furthermore $f \in C^\infty$ then $\partial_i f(x) = 0$ as well.
\end{lemma}

\begin{proof}
\leanok

If $f(x) \ne 0$ then $x$ lies in the support of $f$, hence in $\Omega$,
contradicting $x \notin \Omega$. The same argument applied to $\partial_i f$
(whose support is contained in that of $f$) handles the second statement.
\end{proof}

\begin{theorem}[Zero-extension of a $\Wop{p}$ function]
\label{thm:zero_extend}
\lean{DeGiorgi.zeroExtend_memW1pWitness_p}
\leanok
\uses{def:MemW1pWitness, def:MemW01p}
Let $\Omega \subseteq E$ be open, $1 < p < \infty$, and $v \in \Wop{p}(\Omega)$
with $\Wkp{1}{p}$-witness $G_v$; assume in addition that $v \in W_0^{1,p}(\Omega)$. Define $\tilde v := \mathbf{1}_\Omega \cdot v$
on $E$ and let $\tilde G$ be the vector field whose components are
$\mathbf{1}_\Omega \cdot (G_v)_i$. Then $\tilde G$ is a $\Wkp{1}{p}$-witness
for $\tilde v$ on the whole space $E$.
\end{theorem}

\begin{proof}
\leanok
\uses{lem:weak_partial_unique, lem:weak_of_contDiff, lem:weak_lp_closure, lem:zero_outside_tsupport}

The function $v$ admits an approximating sequence $(\varphi_n)$ of smooth
compactly supported test functions with $\supp\varphi_n \subseteq \Omega$, and
$\varphi_n \to v$, $\partial_i\varphi_n \to (G_v)_i$ in $\Lp{p}(\Omega)$. By
Lemma~\ref{lem:weak_partial_unique} the gradient field $G_v$ obtained from the
arbitrary witness coincides almost everywhere on $\Omega$ with the one from the
$\Wop{p}$ data.

The $\Lp{p}$ membership of $\tilde v$ on $E$ is equivalent to the $\Lp{p}$
membership of $v$ on $\Omega$, via
\leanname{MeasureTheory.memLp\_indicator\_iff\_restrict}, and similarly for each
component of $\tilde G$.

To check that $\tilde G$ is a weak gradient on $E$, we apply
Lemma~\ref{lem:weak_lp_closure} (closure of weak partial derivatives in
$\Lp{p}$, on $\Omega = E$) with the same approximating sequence
$(\varphi_n)$ and the classical gradients $\partial_i\varphi_n$. The
key observations are:
\begin{itemize}
  \item Outside $\Omega$ both $\varphi_n$ and $\partial_i\varphi_n$ vanish (by
        Lemma~\ref{lem:zero_outside_tsupport}); hence
        $\varphi_n - \tilde v = \mathbf{1}_\Omega(\varphi_n - v)$ and
        $\partial_i\varphi_n - \tilde G_i =
         \mathbf{1}_\Omega(\partial_i\varphi_n - (G_v)_i)$ as functions on $E$.
  \item The $\Lp{p}$-norm of an indicator of a measurable set equals the
        $\Lp{p}$-norm computed on the restricted measure
        (\leanname{MeasureTheory.eLpNorm\_indicator\_eq\_eLpNorm\_restrict}),
        so the convergence on $\Omega$ implies convergence on $E$.
\end{itemize}
The hypotheses of Lemma~\ref{lem:weak_lp_closure} are then satisfied, and
$\tilde G_i$ is the weak $i$-th partial derivative of $\tilde v$ on $E$.
\end{proof}

\section{Convolution smoothing}

For $n \in \N$, define a $C^\infty$ bump function on $E$ centered at $0$ with
inner radius $r_{\text{in}}^{(n)} := \frac{1}{2(n+1)}$ and outer radius
$r_{\text{out}}^{(n)} := \frac{1}{n+1}$ (\leanname{DeGiorgi.shrinkingBump} (private)).
We denote the corresponding mass-one normalized bump by $\rho_n :=
(\text{shrinkingBump}_n).\text{normed}$, so $\int_E \rho_n = 1$, $\rho_n \ge 0$,
and $\supp \rho_n \subseteq \Bz{r_{\text{out}}^{(n)}}$. The radii satisfy
\begin{equation*}
  r_{\text{out}}^{(n)} = \frac{1}{n+1} \;\to\; 0
  \quad\text{as } n \to \infty
\end{equation*}
(\leanname{DeGiorgi.tendsto\_shrinkingBump\_rOut} (private)).

\begin{lemma}[Support of a convolution]
\label{lem:tsupport_conv}
\leanok
\noindent\textit{(Formalized as private declarations \leanname{DeGiorgi.tsupport\_convolution\_subset\_cthickening}, \leanname{DeGiorgi.tsupport\_normed\_convolution\_subset\_cthickening}.)}\par
If $\supp \kappa \subseteq \overline{\Bz{r}}$ and $u$ has compact support, then
\begin{equation*}
  \supp(\kappa * u) \;\subseteq\; \{x \in E : \dist(x, \supp u) \le r\}.
\end{equation*}
In particular, $\supp(\rho_n * u) \subseteq \{x : \dist(x,\supp u) \le
r_{\text{out}}^{(n)}\}$.
\end{lemma}

\begin{proof}
\leanok

Pointwise, $(\kappa * u)(x) = \int \kappa(y) u(x - y)\,dy$ vanishes unless
there exists $y \in \supp\kappa$ and $z \in \supp u$ with $x = y + z$, hence
$\dist(x, z) = \|y\| \le r$. Closing under the topology, the closure is
contained in the closed $r$-thickening, which is closed.
\end{proof}

\begin{proposition}[$\Lp{p}$-contraction of normed convolution]
\label{prop:Lp_contraction_conv}
\leanok
\noindent\textit{(Formalized as private declaration \leanname{DeGiorgi.eLpNorm\_normedConvolution\_le}.)}\par
Let $\rho \in C_c^\infty(E)$ with $\rho \ge 0$ and $\int_E \rho = 1$, and let
$1 \le p < \infty$. Then for every $f \in \Lp{p}(E)$,
\begin{equation*}
  \|\rho * f\|_{\Lp{p}(E)} \le \|f\|_{\Lp{p}(E)}.
\end{equation*}
\end{proposition}

\begin{proof}
\leanok

Define the probability measure $\nu$ on $E$ by $d\nu = \rho\,d\,\mathrm{vol}$,
which is a probability measure since $\rho \ge 0$ and $\int \rho = 1$. For each
$x \in E$,
\begin{equation*}
  (\rho * f)(x) = \int_E f(x - y)\,d\nu(y).
\end{equation*}
By the integral form of Jensen's inequality applied to the convex function
$t \mapsto t^p$ on $[0,\infty)$ (\leanname{ConvexOn.map\_integral\_le}),
\begin{equation*}
  |(\rho * f)(x)|^p
    \le \left(\int_E |f(x-y)|\,d\nu(y)\right)^p
    \le \int_E |f(x-y)|^p \,d\nu(y)
    = (\rho * |f|^p)(x).
\end{equation*}
Integrating in $x$ and using Fubini together with $\int \rho = 1$,
\begin{equation*}
  \int_E |(\rho * f)(x)|^p\,dx \le \int_E (\rho * |f|^p)(x)\,dx
    = \left(\int_E \rho\right)\left(\int_E |f|^p\right)
    = \|f\|_{\Lp{p}}^p.
\end{equation*}
Taking $p$-th roots gives the claim.
\end{proof}

\begin{corollary}[Convolution is linear on differences]
\label{cor:conv_sub}
\leanok
\noindent\textit{(Formalized as private declarations \leanname{DeGiorgi.normedConvolution\_sub\_eq}, \leanname{DeGiorgi.eLpNorm\_normedConvolution\_sub\_le}.)}\par
Under the assumptions of Proposition~\ref{prop:Lp_contraction_conv} we also
have, for $f, g \in \Lp{p}(E)$,
\begin{equation*}
  \|\rho * f - \rho * g\|_{\Lp{p}(E)} = \|\rho * (f - g)\|_{\Lp{p}(E)}
   \le \|f - g\|_{\Lp{p}(E)}.
\end{equation*}
\end{corollary}

\begin{proof}
\leanok
\uses{prop:Lp_contraction_conv}

Linearity of $f \mapsto \rho * f$ on locally integrable functions yields the
equality, after which Proposition~\ref{prop:Lp_contraction_conv} applies to
$f - g$.
\end{proof}

\begin{lemma}[Convolution by a shrinking bump approximates a continuous function]
\label{lem:conv_approx_cont}
\leanok
\noindent\textit{(Formalized as private declaration \leanname{DeGiorgi.tendsto\_eLpNorm\_normedConvolution\_sub\_of\_continuous}.)}\par
Let $g \in C_c(E)$ and $1 \le p < \infty$. Then
\begin{equation*}
  \|\rho_n * g - g\|_{\Lp{p}(E)} \;\to\; 0 \quad\text{as } n \to \infty.
\end{equation*}
\end{lemma}

\begin{proof}
\leanok
\uses{lem:zero_outside_tsupport, lem:tsupport_conv}

Choose $S := \{x : \dist(x, \supp g) \le 1\}$ and $T := \{x : \dist(x, \supp g)
\le 2\}$, both compact, with $S \subseteq T$. The function $g$ is uniformly
continuous on the compact set $T$. Fix $\varepsilon > 0$ and let $\varepsilon_0
:= \varepsilon / (\mathrm{vol}(S)^{1/p} + 1)$. By uniform continuity, there
exists $\delta_0 > 0$ such that $|g(y) - g(x)| < \varepsilon_0$ whenever
$x, y \in T$ with $\dist(x, y) < \delta_0$. Set $\delta := \min(\delta_0, 1)$.

For $n$ large enough that $r_{\text{out}}^{(n)} < \delta$, we have $\supp(\rho_n
* g) \subseteq S$ (Lemma~\ref{lem:tsupport_conv}, since
$r_{\text{out}}^{(n)} \le 1$). Pointwise, for $x \in S$,
\begin{equation*}
  |(\rho_n * g)(x) - g(x)|
    = \left|\int_E \rho_n(y) (g(x - y) - g(x))\,dy\right|
    \le \varepsilon_0,
\end{equation*}
since $\rho_n$ is supported in $\Bz{r_{\text{out}}^{(n)}} \subseteq \Bz{\delta}$
and both $x, x - y \in T$ for $y \in \supp\rho_n$. For $x \notin S$ both
$(\rho_n * g)(x)$ and $g(x)$ vanish.

The bound $\|\rho_n * g - g\|_{\Lp{p}(E)} \le \varepsilon_0
\,\mathrm{vol}(S)^{1/p}$ then follows from
\leanname{eLpNorm\_sub\_le\_of\_dist\_bdd}, and the choice of $\varepsilon_0$
gives this is less than $\varepsilon$.
\end{proof}

\begin{theorem}[Convolution by a shrinking bump approximates an $\Lp{p}$ function]
\label{thm:conv_approx_Lp}
\leanok
\noindent\textit{(Formalized as private declaration \leanname{DeGiorgi.tendsto\_eLpNorm\_normedConvolution\_sub}.)}\par
Let $1 \le p < \infty$ and $f \in \Lp{p}(E)$. Then
\begin{equation*}
  \|\rho_n * f - f\|_{\Lp{p}(E)} \;\to\; 0 \quad\text{as } n \to \infty.
\end{equation*}
\end{theorem}

\begin{proof}
\leanok
\uses{cor:conv_sub, lem:conv_approx_cont}

By density of smooth compactly supported functions in $\Lp{p}$ (Mathlib's
\leanname{MemLp.exist\_eLpNorm\_sub\_le}), for any $\varepsilon > 0$ there
exists $g \in C_c^\infty(E)$ with $\|f - g\|_{\Lp{p}} < \varepsilon/3$. Then by
the triangle inequality
\begin{equation*}
  \|\rho_n * f - f\|_{\Lp{p}} \le
  \|\rho_n * f - \rho_n * g\|_{\Lp{p}} + \|\rho_n * g - g\|_{\Lp{p}} +
  \|g - f\|_{\Lp{p}}.
\end{equation*}
By Corollary~\ref{cor:conv_sub} the first term is at most $\|f - g\|_{\Lp{p}} <
\varepsilon/3$. By Lemma~\ref{lem:conv_approx_cont} the middle term tends to
$0$, so for $n$ large it is less than $\varepsilon/3$. The last term is also
less than $\varepsilon/3$. Hence $\|\rho_n * f - f\|_{\Lp{p}} < \varepsilon$
for all sufficiently large $n$, proving convergence to zero.
\end{proof}

\section{Uniqueness of the weak gradient and global smooth approximation}

\begin{theorem}[Uniqueness of the weak gradient]
\label{thm:weak_grad_unique}
\lean{DeGiorgi.MemW1pWitness.ae_eq_p, DeGiorgi.MemW1pWitness.ae_eq}
\leanok
\uses{def:MemW1pWitness}
Let $\Omega \subseteq E$ be open, $1 \le p < \infty$, and $f$ a function with
two $\Wkp{1}{p}$-witnesses on $\Omega$ giving rise to weak gradient fields
$G_1, G_2$. Then $G_1 = G_2$ a.e.\ on $\Omega$. The same holds for $H^1$
witnesses ($p = 2$).
\end{theorem}

\begin{proof}
\leanok
\uses{lem:weak_partial_unique}

Apply Lemma~\ref{lem:weak_partial_unique} componentwise to $(G_1)_i$ and
$(G_2)_i$, then assemble: a.e.\ all components agree, so the vectors
themselves agree a.e.
\end{proof}

\begin{lemma}[Convolution differentiates against the weak partial derivative]
\label{lem:conv_fderiv_weak_partial}
\leanok
\uses{def:HasWeakPartialDeriv}
\noindent\textit{(Formalized as private declaration \leanname{DeGiorgi.convolution\_fderiv\_eq\_convolution\_weakPartial\_univ}.)}\par
Let $u \in L^1_{loc}(E)$ and let $g$ be the weak $i$-th partial derivative of
$u$ on $E$. Then for every $\varphi \in C_c^\infty(E)$ and every $x \in E$,
\begin{equation*}
  \bigl((\partial_i\varphi) * u\bigr)(x) = (\varphi * g)(x).
\end{equation*}
\end{lemma}

\begin{proof}
\leanok

For fixed $x$, define the test function $\psi(t) := \varphi(x - t)$, which is
smooth and compactly supported. A direct computation shows
$(\partial_i \psi)(t) = -(\partial_i\varphi)(x - t)$ (chain rule with the
sign-flipping translation). The weak-derivative identity for $u$ tested against
$\psi$ then reads
\begin{equation*}
  \int_E u(t)\,(\partial_i\psi)(t)\,dt = -\int_E g(t)\,\psi(t)\,dt,
\end{equation*}
i.e.,
\begin{equation*}
  -\int_E u(t)\,(\partial_i\varphi)(x - t)\,dt = -\int_E g(t)\,\varphi(x - t)\,dt.
\end{equation*}
Cancelling the minus signs and rewriting both sides as convolutions yields the
claim.
\end{proof}

\begin{theorem}[Smooth compactly supported approximation on $E$ for compactly
supported $\Wkp{1}{p}$ functions]
\label{thm:smooth_approx_univ}
\lean{DeGiorgi.exists_smooth_compactSupport_W1p_approx_univ}
\leanok
\uses{def:MemW1pWitness}
Let $1 < p < \infty$, and let $u : E \to \R$ be compactly supported with a
$\Wkp{1}{p}$-witness on $E$ with weak gradient $G$. Then there exists a sequence
$(\varphi_n)$ of smooth compactly supported functions on $E$ with
\begin{align*}
  \supp \varphi_n &\subseteq \{x : \dist(x, \supp u) \le r_{\text{out}}^{(n)}\}, \\
  \|\varphi_n - u\|_{\Lp{p}(E)} &\to 0, \\
  \|\partial_i\varphi_n - G_i\|_{\Lp{p}(E)} &\to 0 \quad\text{for every } i.
\end{align*}
\end{theorem}

\begin{proof}
\leanok
\uses{lem:zero_outside_tsupport, lem:tsupport_conv, thm:conv_approx_Lp, lem:conv_fderiv_weak_partial}

Define $\varphi_n := \rho_n * u$. Since $u$ is compactly supported and $\rho_n
\in C_c^\infty(E)$, $\varphi_n$ is smooth and has support contained in the
$r_{\text{out}}^{(n)}$-thickening of $\supp u$ (Lemma~\ref{lem:tsupport_conv}).
The $\Lp{p}$ approximation $\varphi_n \to u$ is exactly
Theorem~\ref{thm:conv_approx_Lp} applied to $u$.

For the gradients, the classical fact
\begin{equation*}
  \partial_i (\rho_n * u) = (\partial_i \rho_n) * u
\end{equation*}
(Mathlib's
\leanname{HasCompactSupport.hasFDerivAt\_convolution\_left})
combined with Lemma~\ref{lem:conv_fderiv_weak_partial} (with $u$ on the right
and the smooth $\rho_n$ on the left) yields
\begin{equation*}
  \partial_i \varphi_n = (\partial_i\rho_n) * u = \rho_n * G_i,
\end{equation*}
since $G_i$ is the weak partial derivative of $u$ on $E$. Applying
Theorem~\ref{thm:conv_approx_Lp} to $G_i \in \Lp{p}(E)$ then gives
$\partial_i\varphi_n - G_i = \rho_n * G_i - G_i \to 0$ in $\Lp{p}(E)$.
\end{proof}

\begin{theorem}[Smooth approximation for a witness supported in a compact subset]
\label{thm:smooth_approx_local}
\lean{DeGiorgi.exists_smooth_W1p_approx_of_supportedWitness}
\leanok
\uses{def:MemW1pWitness}
Let $\Omega$ be open, $1 < p < \infty$, and let $u$ have a $\Wkp{1}{p}$-witness
on $\Omega$ with gradient field $G$. Suppose $K \subseteq \Omega$ is compact
with $\supp u \subseteq K$ and $\supp G_i \subseteq K$ for every $i$. Then
there exists a sequence $(\varphi_n)$ of smooth compactly supported functions
with $\supp\varphi_n \subseteq \Omega$ such that
\begin{align*}
  \|\varphi_n - u\|_{\Lp{p}(\Omega)} &\to 0, \\
  \|\partial_i\varphi_n - G_i\|_{\Lp{p}(\Omega)} &\to 0 \quad\text{for every } i.
\end{align*}
\end{theorem}

\begin{proof}
\leanok
\uses{lem:smooth_cutoff, lem:zero_outside_tsupport, thm:smooth_approx_univ}

Choose $\delta > 0$ so that the closed $\delta$-thickening $K' := \{x :
\dist(x,K) \le \delta\}$ is contained in $\Omega$
(\leanname{IsCompact.exists\_cthickening\_subset\_open}), and a smooth cutoff
$\chi$ as in Lemma~\ref{lem:smooth_cutoff} satisfying $\chi = 1$ on $K'$ and
$\supp \chi \subseteq \Omega$.

We promote the local witness on $\Omega$ to a witness on the whole space $E$
with the same gradient field. The $\Lp{p}$-membership extends to all of $E$
because $u$ and each $G_i$ vanish outside $K \subseteq \Omega$
(Lemma~\ref{lem:zero_outside_tsupport}). To verify the weak-gradient identity
on $E$, take any test function $\varphi \in C_c^\infty(E)$ and consider $\psi
:= \chi \varphi$. Then $\psi \in C_c^\infty$ and $\supp \psi \subseteq \Omega$,
so the witness on $\Omega$ gives
\begin{equation*}
  \int_\Omega u(x)\,\partial_i\psi(x)\,dx = -\int_\Omega G_i(x)\,\psi(x)\,dx.
\end{equation*}
On $K$ we have $\chi = 1$ and $\partial_i\chi = 0$ (because $\chi$ is constant
equal to $1$ on the larger neighborhood $K'$ of $K$); since $u$ and $G_i$ are
both supported in $K$, this gives pointwise on $E$
\begin{equation*}
  u(x)\,\partial_i\psi(x) = u(x)\,\partial_i\varphi(x), \qquad
  G_i(x)\,\psi(x) = G_i(x)\,\varphi(x),
\end{equation*}
and the integral identity transfers from $\psi$ to $\varphi$.

Now apply Theorem~\ref{thm:smooth_approx_univ} on the whole space to get a
sequence $(\widetilde\varphi_m)$ with
\begin{equation*}
  \supp \widetilde\varphi_m \subseteq \{x : \dist(x, \supp u) \le r_{\text{out}}^{(m)}\}.
\end{equation*}
For $m$ large enough that $r_{\text{out}}^{(m)} \le \delta$, this support is
contained in $K' \subseteq \Omega$. Setting $\varphi_n :=
\widetilde\varphi_{n+N}$ for an appropriate $N$ furnishes the desired sequence:
the support condition holds for every $n$, and $\Lp{p}(\Omega)$-convergence
follows from $\Lp{p}(E)$-convergence by monotonicity of the $\Lp{p}$-norm in
the underlying measure (\leanname{eLpNorm\_mono\_measure}).
\end{proof}

\begin{theorem}[Localization: compactly supported $\Wkp{1}{p}$ implies $\Wop{p}$]
\label{thm:memW01p_of_compactly_supported}
\lean{DeGiorgi.memW01p_of_memW1p_of_tsupport_subset}
\leanok
\uses{def:MemW1p, def:MemW01p}
Let $\Omega \subseteq E$ be open, $1 < p < \infty$, and $u \in W^{1,p}(\Omega)$
with compact support and $\supp u \subseteq \Omega$. Then $u \in \Wop{p}(\Omega)$.
\end{theorem}

\begin{proof}
\leanok
\uses{def:MemW1pWitness, prop:witness_mul_smooth_bounded, lem:smooth_cutoff, thm:smooth_approx_local}

Choose any $\Wkp{1}{p}$-witness for $u$ on $\Omega$, with gradient field $G$.
Pick $\delta > 0$ such that the closed $\delta$-thickening of $\supp u$ is
contained in $\Omega$, and let $K$ be that thickening. Apply
Lemma~\ref{lem:smooth_cutoff} to obtain a smooth cutoff $\eta$ equal to $1$ on
$K$ with $\supp \eta \subseteq \Omega$. Since $\eta$ is smooth and compactly
supported, both $|\eta|$ and $\|D\eta\|$ are bounded by some constants $C_0 :=
1$ and $C_1 \ge 0$ (the latter using continuity of $D\eta$ on its compact
support, \leanname{HasCompactSupport.exists\_bound\_of\_continuousOn}).

By Proposition~\ref{prop:witness_mul_smooth_bounded} the function $v := \eta u$
admits a $\Wkp{1}{p}$-witness. Since $\eta = 1$ on a neighborhood of $\supp u$
and $u = 0$ outside $\supp u$, we have $v = u$ pointwise. Moreover,
$\supp v \subseteq \supp \eta$ (a compact subset of $\Omega$), and the same is
true for each component of the new gradient, by
\leanname{DeGiorgi.tsupport\_mul\_smooth\_bounded\_p\_weakGrad\_component\_subset}
(private).

We are now in the situation of Theorem~\ref{thm:smooth_approx_local} with $K :=
\supp\eta$. The resulting smooth approximating sequence shows $u = v$ belongs
to $\Wop{p}(\Omega)$.
\end{proof}

\section{Whole-space Sobolev inequality}

\begin{definition}[Sobolev constant]
\label{def:C_gns}
\lean{DeGiorgi.C_gns}
\leanok
For $d \in \N_{\ge 1}$ and $p \in \R$, define
\begin{equation*}
  C_{\mathrm{gns}}(d, p) :=
  \mathrm{SNormLESNormFDerivOfEqConst_{\R}}\bigl(\mathrm{vol}_{\R^d}, p\bigr),
\end{equation*}
the constant supplied by Mathlib's general $W^{1,p} \hookrightarrow L^{p^*}$
embedding. By construction, $C_{\mathrm{gns}}(d,p) \ge 0$
(\leanname{DeGiorgi.C\_gns\_nonneg}).
\end{definition}

\begin{lemma}[Sobolev exponent]
\label{lem:sobolev_exp}
\leanok
\noindent\textit{(Formalized as private declarations \leanname{DeGiorgi.sobolev\_exp\_pos}, \leanname{DeGiorgi.sobolev\_exponent\_relation}.)}\par
Let $1 \le p < d$. Then $\frac{dp}{d - p} > 0$, and
\begin{equation*}
  \left(\frac{dp}{d - p}\right)^{-1} = \frac{1}{p} - \frac{1}{d}.
\end{equation*}
\end{lemma}

\begin{proof}
\leanok

Both numerator $dp$ and denominator $d - p$ are positive (since $d \ge 1$ and
$p < d$). The exponent identity follows by direct algebraic manipulation
$\frac{d - p}{dp} = \frac{1}{p} - \frac{1}{d}$.
\end{proof}

\begin{theorem}[Sobolev inequality for smooth compactly supported functions]
\label{thm:sobolev_smooth}
\lean{DeGiorgi.sobolev_smooth}
\leanok
\uses{def:C_gns}
Let $1 \le p < d$, $u \in C_c^1(E)$. Then
\begin{equation*}
  \|u\|_{L^{p^*}(E)} \le C_{\mathrm{gns}}(d, p)\,\|Du\|_{L^p(E)},
  \qquad p^* = \frac{dp}{d-p}.
\end{equation*}
\end{theorem}

\begin{proof}
\leanok
\uses{lem:sobolev_exp}

This is a direct application of Mathlib's general result
\leanname{MeasureTheory.eLpNorm\_le\_eLpNorm\_fderiv\_of\_eq} with the exponent
relation supplied by Lemma~\ref{lem:sobolev_exp} (rewritten in terms of
$\mathrm{Real.toNNReal}$ to match the Mathlib API).
\end{proof}

\begin{lemma}[Operator norm of the differential equals partials norm]
\label{lem:norm_fderiv}
\leanok
\noindent\textit{(Formalized as private declaration \leanname{DeGiorgi.norm\_fderiv\_eq\_norm\_partials}.)}\par
For any $f : E \to \R$ with derivative $Df(x)$ at $x$,
\begin{equation*}
  \|Df(x)\|
   = \left(\sum_{i=1}^d |\partial_i f(x)|^2\right)^{1/2}.
\end{equation*}
\end{lemma}

\begin{proof}
\leanok

By Riesz representation in the Hilbert space $E$, $Df(x)$ corresponds to a
unique $v \in E$ with $Df(x) = \langle v, \cdot\rangle$, and $\|Df(x)\| = \|v\|$.
Evaluating at the standard basis vector $e_i$ gives $v_i = Df(x)(e_i) =
\partial_i f(x)$, so $\|v\|^2 = \sum_i |\partial_i f(x)|^2$.
\end{proof}

\begin{theorem}[Sobolev inequality from approximating sequences]
\label{thm:sobolev_of_approx}
\lean{DeGiorgi.sobolev_of_approx}
\leanok
\uses{def:C_gns}
Let $1 \le p < d$, $u : E \to \R$ a.e.\ strongly measurable, and $G : E \to E$
a vector field whose components are a.e.\ strongly measurable. Suppose there
is a sequence $(\varphi_n)$ of smooth compactly supported functions on $E$
with
\begin{equation*}
  \|\varphi_n - u\|_{\Lp{p}(E)} \to 0, \qquad
  \|\partial_i\varphi_n - G_i\|_{\Lp{p}(E)} \to 0 \;\text{ for every } i.
\end{equation*}
Then
\begin{equation*}
  \|u\|_{L^{p^*}(E)} \le C_{\mathrm{gns}}(d,p)\,\|\,\|G(\cdot)\|\,\|_{\Lp{p}(E)},
  \qquad p^* = \frac{dp}{d-p}.
\end{equation*}
\end{theorem}

\begin{proof}
\leanok
\uses{thm:sobolev_smooth, lem:norm_fderiv}

Apply Theorem~\ref{thm:sobolev_smooth} to each $\varphi_n$:
\begin{equation*}
  \|\varphi_n\|_{L^{p^*}} \le C_{\mathrm{gns}}(d,p)\,\|D\varphi_n\|_{L^p}.
\end{equation*}
Define $\Phi_n(x) := (\partial_1\varphi_n(x),\ldots,\partial_d\varphi_n(x))$,
the vector of classical partial derivatives. By Lemma~\ref{lem:norm_fderiv},
$\|D\varphi_n(x)\| = \|\Phi_n(x)\|$, hence $\|D\varphi_n\|_{L^p} = \|\,\|\Phi_n\|\,\|_{L^p}$.

By the triangle inequality applied componentwise (and bounding the Euclidean
norm of a vector by the sum of the norms of its components, see
\leanname{DeGiorgi.gradVec\_eLpNorm\_le\_sum} (private)),
\begin{equation*}
  \|\Phi_n - G\|_{L^p(E; E)}
   \le \sum_{i=1}^d \|\partial_i\varphi_n - G_i\|_{\Lp{p}(E)} \;\to\; 0.
\end{equation*}
Hence $\|\,\|\Phi_n\|\,\|_{L^p} \to \|\,\|G\|\,\|_{L^p}$, by reverse triangle
inequality $|\,\|\Phi_n(x)\| - \|G(x)\|\,| \le \|\Phi_n(x) - G(x)\|$.

Convergence in $L^p$ implies convergence in measure
(\leanname{tendstoInMeasure\_of\_tendsto\_eLpNorm}), so by passing to a
subsequence we may assume $\varphi_{\sigma(n)} \to u$ a.e. By Fatou's lemma
applied to $L^{p^*}$ (\leanname{Lp.eLpNorm\_lim\_le\_liminf\_eLpNorm}),
\begin{equation*}
  \|u\|_{L^{p^*}}
   \le \liminf_{n\to\infty} \|\varphi_{\sigma(n)}\|_{L^{p^*}}
   \le \liminf_{n\to\infty} C_{\mathrm{gns}}(d,p)\,\|\,\|\Phi_{\sigma(n)}\|\,\|_{L^p}.
\end{equation*}
The right-hand side $\liminf$ is bounded above by $C_{\mathrm{gns}}(d,p)\,\|
\,\|G\|\,\|_{L^p}$, since $\|\,\|\Phi_{\sigma(n)}\|\,\|_{L^p} \le \|\,\|G\|\,\|_{L^p}
+ \|\Phi_{\sigma(n)} - G\|_{L^p}$ and the second term tends to zero, using
\leanname{ENNReal.liminf\_add\_of\_right\_tendsto\_zero}. Combining yields the
claim.
\end{proof}

\begin{corollary}[Sobolev inequality on $\Wop{p}(E)$]
\label{cor:sobolev_W01p}
\lean{DeGiorgi.sobolev_of_memW01p_univ}
\leanok
\uses{def:MemW1pWitness, def:MemW01p, def:C_gns}
Let $1 \le p < d$ and $u \in \Wop{p}(E)$. Then there exists a
$\Wkp{1}{p}$-witness for $u$ on $E$ such that
\begin{equation*}
  \|u\|_{L^{p^*}(E)} \le C_{\mathrm{gns}}(d,p)\,\|\,\|G(\cdot)\|\,\|_{\Lp{p}(E)},
  \qquad p^* = \frac{dp}{d - p},
\end{equation*}
where $G$ is the weak gradient field carried by the witness.
\end{corollary}

\begin{proof}
\leanok
\uses{thm:sobolev_of_approx}

By Definition~\ref{def:MemW01p}, $u \in \Wop{p}(E)$ comes equipped with a
$\Wkp{1}{p}$-witness on $E$ and an approximating sequence $(\varphi_n)$ of
smooth compactly supported functions converging to $u$ and whose classical
gradients converge to the weak gradient $G$ in $\Lp{p}$. The hypotheses of
Theorem~\ref{thm:sobolev_of_approx} are satisfied, and its conclusion is
exactly the desired inequality.
\end{proof}

\section{Positive-part prelude}

\begin{lemma}[Indicator truncation preserves $L^2$]
\label{lem:indicator_component_memLp}
\lean{DeGiorgi.indicator_component_memLp}
\leanok
Let $\Omega \subseteq E$, $\sigma : E \to \R$ be a.e.\ measurable on $\Omega$
and $g \in \Lp{2}(\Omega)$. Then the function
\begin{equation*}
  x \mapsto \mathbf{1}_{\{\sigma > 0\}}(x)\,g(x)
\end{equation*}
also lies in $\Lp{2}(\Omega)$.
\end{lemma}

\begin{proof}
\leanok

The set $\{\sigma > 0\}$ is measurable since $\sigma$ is a.e.\ measurable, so
the indicator multiplication preserves a.e.\ strong measurability. Pointwise,
$|\mathbf{1}_{\{\sigma > 0\}}(x)\,g(x)| \le |g(x)|$, and squaring, integrating
and applying the $\Lp{2}$ membership of $g$ gives the conclusion.
\end{proof}

\begin{lemma}[Truncated weak-gradient component in $L^2$]
\label{lem:posPartGrad}
\lean{DeGiorgi.posPartGrad_component_memLp}
\leanok
\uses{def:MemW1pWitness}
Let $u \in H^1(\Omega)$ with $H^1$-witness gradient field $G$. For each
$i \in \{1,\ldots,d\}$, the function
\begin{equation*}
  x \mapsto \mathbf{1}_{\{u > 0\}}(x)\,G_i(x)
\end{equation*}
belongs to $\Lp{2}(\Omega)$.
\end{lemma}

\begin{proof}
\leanok
\uses{lem:indicator_component_memLp}

Apply Lemma~\ref{lem:indicator_component_memLp} with $\sigma := u$ and $g :=
G_i$, both of which lie in the relevant $\Lp{2}$ spaces (the second by the
witness assumption).
\end{proof}


%% ========================================================================
